\documentclass{article}
\usepackage[T1]{fontenc}
\usepackage{lmodern}
\usepackage{graphicx}
\usepackage{amsmath,amssymb}
\usepackage[a4paper,margin=2cm]{geometry}
\usepackage[absolute,overlay]{textpos}
\setlength{\TPHorizModule}{1mm}
\setlength{\TPVertModule}{1mm}
\usepackage[indonesian]{babel}
\usepackage{hyperref}
\setlength{\emergencystretch}{2em}
% unit: Halaman 74

\begin{document}

% === Halaman 74 (llm) ===

\section*{3. Trivium}

\begin{itemize}
    \item Trivium adalah \textit{cipher-alir} modern, merupakan salah satu \textit{cipher} alir yang di-submit pada kompetisi \textit{stream cipher} eSTREAM pada tahun 2004 -- 2008.
    \item Dibuat oleh Christophe De Cannière dan Bart Preneel.
    \item Menggunakan tiga buah \textit{nonlinear LFSRs} (NLFSR) dengan panjang masing-masing 93, 84, dan 111 bit.
    \item Melakukan XOR terhadap tiga buah luaran NLFSR untuk membangkitkan bit kunci-alir.
    \item Minimalis jika diimplementasikan pada \textit{hardware}:
    \begin{itemize}
        \item total sel register: 288
        \item tiga buah gerbang AND
        \item tujuh buah gerbang XOR
    \end{itemize}
\end{itemize}

\end{document}
